\documentclass{article}
\usepackage{hyperref}
\usepackage[a4paper]{geometry}
\usepackage{fancyhdr}
\pagestyle{fancy}
\lhead{Elektronenbeugung}
\rhead{März 2026}
\begin{document}
\section{Elektronenbeugung}
Aus einer Elektronenkanone werden Elektronen mit einer bestimmten, konstanten Geschwindigkeit geschossen. Diese treffen auf polykristalline Graphitplättchen.
 
Dabei erzeugen sie auf dem nachfolgenden Leuchtschirm ein radiales Interferenzmuster aus unterschiedlich großen, zentrierten Ringen.
 
\subsection{Wie das Gitter}
Visuell allein fällt auf, dass das erzeugte Muster dem Interferenzmuster des \hyperref[Das Gitter]{Gitters}, aufgenommen mit Licht, ähnelt. 
 
Das Interferenzmuster des Lichtes am Gitter basierte aber auf den Welleneigenschaften des Lichtes, so folgt also, dass Elektronen auch Welleneigenschaften brauchen. So haben auch Elektronen und andere Quanten eine Wellenlänge, die \hyperref[Die de-Broglie-Wellenlänge]{de-Broglie-Wellenlänge}. 
\end{document}